\begin{claim}{D.2.39}{Proved}{Classifying-space smash preserves every reversible class below its own}
Let $p$ be a prime, let $G=B\mathbb Z/p$, and let $C$ be a
connected pointed CW space which is ordinarily $G$-acyclic.
Put
\[
       M=M(\mathbb Z/p,1)=\cofib(p:S^1\to S^1),
       \qquad X=G\wedge C.
\]
Then
\[
 \begin{gathered}
 P_X\bigl(\Sigma^2(M\wedge C)\bigr)\simeq *,\qquad
 P_X(\Sigma^4C)\simeq *,\\
 \min\mathrm{Rev}\{G\wedge C\to *\}
   =\min\mathrm{Rev}\{M\wedge C\to *\}
   =\min\mathrm{Rev}\{C\to *\}.
 \end{gathered}
 \tag{D.2.39a}
\]
More generally, for every $C$ in the aspherical class of
$\min\mathrm{Rev}(G\to *)$, one has
\[
 \min\mathrm{Rev}\{G\wedge C\to *\}
             =\min\mathrm{Rev}\{C\to *\}.
 \tag{D.2.39a'}
\]
Thus smashing the classifying-space/Moore-space distinction with
any such factor cannot retain that distinction between the
generated reversible classes.
No finiteness assumption on $C$, or compactness assumption on its
suspension spectrum, is required.
\end{claim}
\begin{proof}
We first establish the spectral consequence of
$B\mathbb Z/p$-nullity needed here:
\[
 \begin{gathered}
 K\simeq\Omega^\infty E\text{ for a connective spectrum }E,
 \quad\Map_*(G,K)\simeq *\quad\Longrightarrow\\
 p:\pi_jK\xrightarrow{\cong}\pi_jK
                  \quad\text{for every }j\geq2.
 \end{gathered}
 \tag{D.2.39b}
\]
Set $W=\Sigma^\infty G$ and $W_i=\Sigma^{-i}W$ for $i\geq0$.
The adjunction between suspension spectra and infinite-loop
spaces identifies the hypothesis with spectral $W$-nullity of
$E$, using mapping spaces, so $P_W^{\Sp}E\simeq E$.
The spectrum $W$ is $1$-connective and satisfies the
$1$-supported $p$-torsion condition. Write
\[
       E_i=P_{W_i}^{\Sp}E,
       \qquad E_\infty=\colim_{i\geq0}E_i.
 \tag{D.2.39c}
\]
The natural maps point from $E_i$ to $E_{i+1}$.
Bousfield's exact comparison theorem gives
\[
       \fib(E_i\to E_{i+1})\simeq\Sigma^{-i}HG_i
 \tag{D.2.39d}
\]
for some $p$-primary torsion group $G_i$
\cite[Theorem~2.6, p.~37]{Bousfield1996ULP}.
By extension closure of homotopy-group bounds and exactness of
filtered colimits in spectra, this implies
\[
       \fib(E_i\to E_\infty)\in\Ho(\Sp)_{\leq-i}.
 \tag{D.2.39e}
\]

Here is the required check at the limit. Fix $i\geq0$ and use
$E_{i+1}\to E_\infty$. Its fiber lies in degrees at most
$-i-1$, and its suspended fiber lies in degrees at most $-i$.
On the other hand $W_i\in\Ho(\Sp)_{\geq1-i}$.
The two $t$-structure orthogonalities therefore give an
equivalence of entire mapping spaces
\[
 \Map_{\Sp}(W_i,E_{i+1})
          \xrightarrow{\simeq}\Map_{\Sp}(W_i,E_\infty).
 \tag{D.2.39f}
\]
The left side is contractible, since $E_{i+1}$ is
$W_{i+1}$-null and $W_i=\Sigma W_{i+1}$.
It follows that $E_\infty$ is right orthogonal to every
integer suspension of $W$. The established Kahn--Priddy
calculation, C.3 and its prime-$p$ version in D.2.2, gives
\[
             \loc(W)=\loc(\mathbb S/p).
 \tag{D.2.39g}
\]
Consequently the function spectrum
$F(\mathbb S/p,E_\infty)$ is zero, and multiplication by $p$
on $E_\infty$ is an equivalence. Since $E_0=E$, the case
$i=0$ of \textup{(D.2.39e)} identifies $\pi_jE$ with
$\pi_jE_\infty$ for $j\geq2$. This proves
\textup{(D.2.39b)}. In particular, this argument has not
commuted $\Map_{\Sp}(W,-)$ with a filtered colimit; the
comparison with $E_{i+1}$ controls components as well as
positive homotopy groups.

Now let $Z$ be $X$-null. Then $\Map_*(C,Z)$ is $G$-null.
The assumed ordinary $G$-acyclicity of $C$ gives
\[
 \Map_*(C\wedge C,Z)
       \simeq\Map_*(C,\Map_*(C,Z))\simeq *.
 \tag{D.2.39h}
\]
As in D.2.37, all higher smash powers of $C$ are $X$-acyclic,
and the finite product filtration proves that
$C^r/\bigvee_{i=1}^rC$ is $X$-acyclic for every finite $r$.
Replace $\Omega Z$ by a simplicial group $H$. Restriction
\[
       \Map_*(C^r,H)\longrightarrow\Map_*(C,H)^r
 \tag{D.2.39i}
\]
is a Kan fibration of simplicial groups with contractible
zero fiber. Every tuple extends by ordered multiplication
on its separate coordinates. Translation therefore identifies
every fiber with the zero fiber, proving that
\textup{(D.2.39i)} is a weak equivalence on all components.

For this general $C$ there is no linear algebraic theory.
There is instead an actual functor on finite pointed sets:
\[
 \mathcal F(S)=\Map_*(C^{S\setminus\{*\}},H).
 \tag{D.2.39j}
\]
A pointed map $f:S\to T$ pulls coordinates back to give
$C^{T\setminus\{*\}}\to C^{S\setminus\{*\}}$, with the
basepoint in coordinates mapped to $*$. Precomposition gives
the required covariant map $\mathcal F(S)\to\mathcal F(T)$.
This is a $\Gamma$-space in the convention that
$\Gamma^{\mathrm{op}}$ is the category of finite pointed
sets. Its Segal maps are precisely
\textup{(D.2.39i)}, so it is special. The multiplication on
$\pi_0\mathcal F(\{*,1\})$ is induced by the fold map
$\{*,1,2\}\to\{*,1\}$. On powers of $C$ this is the
diagonal, and, using the ordered-product extension in
\textup{(D.2.39i)}, the induced law is pointwise
multiplication in $H$. Its components therefore form a
group. Thus $\mathcal F$ is very special.

Segal's recognition theorem now supplies a connective
spectrum $E$ and a pointed weak equivalence
\[
             K=\Map_*(C,\Omega Z)\simeq\Omega^\infty E
 \tag{D.2.39k}
\]
\cite[Definitions~1.1--1.3 and Proposition~1.4,
pp.~295--296; proof in Propositions~1.5--1.6 and
Lemmas~1.7--1.8, pp.~296--298]{I_Segal1974}.
One can realize the simplicial mapping spaces before applying
that formulation; their weak equivalences are homotopy
equivalences of CW types. Repeated classifying spaces give
the connective spectrum. This is an infinite-loop-space
conclusion, with no assertion that $K$ is a GEM.

By adjunction,
\[
             \Map_*(G,K)
                 \simeq\Map_*(X,\Omega Z)\simeq *.
 \tag{D.2.39l}
\]
Apply \textup{(D.2.39b)}. Multiplication by $p$ is an
isomorphism on all homotopy groups of $\Omega^2K$, including
its group of components. The $p$-power map on $\Omega^2K$,
induced by the degree-$p$ circle map, is consequently an
equivalence. The defining
Moore cofibration gives
\[
 \begin{aligned}
 \Map_*(M,\Omega K)
    &\simeq\fib\bigl(p:\Omega^2K\to\Omega^2K\bigr)
       \simeq *,\\
 \Map_*(M,\Omega K)
    &\simeq\Map_*\bigl(\Sigma^2(M\wedge C),Z\bigr).
 \end{aligned}
 \tag{D.2.39m}
\]
This proves the first ordinary acyclicity in
\textup{(D.2.39a)} and, by desuspension, puts $M\wedge C$
in the reversible null class generated by $X$.

We finish with an elementary telescope argument which also
recovers $C$ itself. More generally, let $D$ be any connected
space with $p$-primary reduced integral homology, put
$T=\Sigma D$, and take the suspension degree map
\[
 f=(p:S^1\to S^1)\wedge\id_D:T\longrightarrow T.
 \tag{D.2.39n}
\]
It induces multiplication by $p$ on reduced integral homology.
Its mapping telescope is simply connected and has zero reduced
integral homology, since each homology element is killed by a
power of $p$. The telescope is consequently contractible.
Set $D_r=\cofib(f^r)$, for $r\geq1$. The cofiber-of-composition
pushout gives actual homotopy cofiber sequences
\[
            D_r\longrightarrow D_{r+1}\longrightarrow D_1,
            \qquad D_1=M\wedge D.
 \tag{D.2.39o}
\]
Thus every $D_r$ is ordinarily $D_1$-acyclic. The structure
maps in this sequence use the identity on the source copy of
$T$ and $f$ on the target copy. Homotopy colimits preserve
these cofibers, so
\[
 \hocolim_rD_r
     \simeq\cofib\bigl(T\to\operatorname{tel}(f)\bigr)
     \simeq\Sigma T=\Sigma^2D.
 \tag{D.2.39p}
\]
In particular $\Sigma^2D$ is ordinarily $(M\wedge D)$-acyclic,
and
\[
       \min\mathrm{Rev}\{M\wedge D\to *\}
                  =\min\mathrm{Rev}\{D\to *\}.
 \tag{D.2.39q}
\]
The other inclusion in \textup{(D.2.39q)} is just closure
under smash products.

An ordinarily $G$-acyclic space has $p$-primary reduced integral
homology, as proved in D.2.7: every $H\mathbb Z[1/p]$-local
space is $G$-null. Apply the preceding construction with
$D=C$, and suspend every cofiber sequence and its final
homotopy colimit twice. Its building space
$\Sigma^2(M\wedge C)$ is $X$-acyclic by
\textup{(D.2.39m)}. Hence its resulting space $\Sigma^4C$
is $X$-acyclic. Desuspension puts $C$ in the reversible
class generated by $G\wedge C$, while $G\wedge C$ belongs
to the class generated by $C$ by smash closure. Together
with \textup{(D.2.39q)}, this proves all of
\textup{(D.2.39a)}.

Finally, if $C$ lies in the reversible null class generated
by $G$, D.2.7 with $r=1$ says that $\Sigma C$ is ordinarily
$G$-acyclic. In particular $C$ has $p$-primary reduced
homology and is connected. Apply the proved case to
$\Sigma C$. Since $G\wedge\Sigma C\simeq\Sigma(G\wedge C)$,
and a space and its suspension generate the same reversible
class, the resulting equality is exactly
\textup{(D.2.39a')}.
\end{proof}
